\documentclass[11pt,a4paper]{article}

% Paquetes necesarios
\usepackage[utf8]{inputenc}
\usepackage[spanish]{babel}
\usepackage[margin=2cm]{geometry}
\usepackage{enumitem}
\usepackage{titlesec}
\usepackage{hyperref}
\usepackage{parskip}
\usepackage{xcolor}

% Configuración de colores
\definecolor{darkblue}{RGB}{0,51,102}

% Configuración de hipervínculos
\hypersetup{
    colorlinks=true,
    linkcolor=darkblue,
    urlcolor=darkblue,
    filecolor=darkblue
}

% Configuración de títulos de sección
\titleformat{\section}
{\large\bfseries\color{darkblue}}
{}
{0em}
{}[\titlerule]

\titlespacing*{\section}{0pt}{10pt}{5pt}

\titleformat{\subsection}
{\bfseries}
{}
{0em}
{}

\titlespacing*{\subsection}{0pt}{6pt}{3pt}

% Configuración de listas
\setlist[itemize]{leftmargin=*,nosep,topsep=2pt}

% Eliminar numeración de páginas
\pagestyle{empty}

% Reducir espacio entre párrafos
\setlength{\parskip}{3pt}

\begin{document}

% ==================== ENCABEZADO ====================
\begin{center}
{\LARGE \textbf{Irvin Benitez}} \\[4pt]
\textit{Desarrollador de Software} \\[6pt]

Panamá Oeste, Panamá \hfill
+507 6361-5832 \hfill
\href{mailto:Irvin.benitezs.26@gmail.com}{Irvin.benitezs.26@gmail.com} \\
\href{https://www.linkedin.com/in/irvin-benitez-11313231b/}{LinkedIn} \hfill
\href{https://irvin-benitez.software}{Portafolio} \\
\end{center}



\vspace{-8pt}

% ==================== PERFIL PROFESIONAL ====================
\section*{PERFIL PROFESIONAL}
Desarrollador de software con experiencia en React Native para desarrollo móvil multiplataforma y front-end en React, Next.js, HTML, CSS, JavaScript y TypeScript, utilizando Git para control de versiones. Implemento flujos CI/CD para despliegue continuo y seguro, aplicando prácticas de DevSecOps. Desarrollo back-end robusto con Django y Flask, con gestión eficiente de bases de datos como SQL Server, SQLite, MySQL, MariaDB y PostgreSQL.

Habilidad en el análisis de datos con pandas y creación de notebooks interactivos utilizando Jupyter, facilitando la visualización y manipulación de información para apoyar la toma de decisiones. Integro principios de diseño como UI/UX, MVC y UML para crear sistemas alineados con la lógica del negocio y las necesidades del usuario.

En proyectos individuales y colaborativos utilizo metodologías ágiles como Scrum (sprints, retrospectivas) para entregar software escalable, mantenible y centrado en el usuario.

% ==================== EXPERIENCIA LABORAL ====================
\section*{EXPERIENCIA LABORAL}

\subsection*{Ayudante Académico Estudiantil -- Desarrollador Full-Stack}
\textit{Universidad Tecnológica de Panamá (Alpha Persei)}\\
\textit{Abril 2025 -- Presente}

Colaboro en el desarrollo de un sistema avanzado de información científica para la Universidad Tecnológica de Panamá (UTP), utilizando Django con templates y PostgreSQL como base de datos. Mis contribuciones específicas incluyen la implementación de funcionalidades para la centralización de datos de investigación institucional, integración automatizada con sistemas externos como ORCID para sincronización de perfiles académicos, y desarrollo del módulo de generación automatizada de informes según estándares internacionales (CONACYT, SENACYT). Este trabajo ha optimizado significativamente los procesos de reportería y gestión de información científica de la universidad.

\vspace{4pt}

\subsection*{Pasante -- Desarrollador Full-Stack}
\textit{Universidad Tecnológica de Panamá (DEMI)}\\
\textit{Febrero -- Marzo 2025}

Desarrollé una plataforma web completa para la gestión de mantenimiento institucional en la Universidad Tecnológica de Panamá (UTP), utilizando Flask (Python), Vue.js, y MariaDB. La plataforma permite la gestión integral de solicitudes, programación de mantenimientos preventivos y correctivos, asignación de tareas, y control de inventario. Implementé un sistema de autenticación con credenciales institucionales, un tablero de seguimiento en tiempo real y un calendario interactivo accesible desde cualquier dispositivo conectado al campus universitario, mejorando la eficiencia operativa del departamento.

% ==================== PROYECTOS DESTACADOS ====================
\section*{PROYECTOS DESTACADOS}

\subsection*{Chatbot de WhatsApp Impulsado con Inteligencia Artificial}
Desarrollo de un bot personalizado de WhatsApp con respuestas inteligentes y adaptadas, utilizando la API de Gemini AI. Integración de funciones avanzadas para ofrecer atención al cliente eficiente y personalizada, con capacidad de procesamiento de lenguaje natural y respuestas contextuales.

\vspace{4pt}

\subsection*{Página Web para Empresa de Diseño Gráfico}
Creación de una página web completa para una empresa de diseño gráfico, utilizando PHP y MySQL. Incluye funcionalidades como sistema de inicio de sesión y registro de usuarios, gestión de servicios y portafolio, y un panel de administración completo para manejar contenido, usuarios y servicios ofrecidos.

% ==================== EDUCACIÓN ====================
\section*{EDUCACIÓN}

\subsection*{Licenciatura en Desarrollo y Gestión de Software}
\textit{Universidad Tecnológica de Panamá}\\
\textit{Marzo 2023 -- Diciembre 2026}

% ==================== HABILIDADES TÉCNICAS ====================
\section*{HABILIDADES TÉCNICAS}

\textbf{Frontend:} React, Next.js, Vue.js, React Native, HTML, CSS, JavaScript, TypeScript, Tailwind

\textbf{Backend:} Django, Flask, Python, PHP, Laravel

\textbf{Bases de Datos:} PostgreSQL, MySQL, MariaDB, SQLite, SQL Server

\textbf{Herramientas:} Git/GitHub, Docker, CI/CD, Jupyter Notebook

\textbf{Análisis de Datos:} Pandas, Data visualization

\textbf{Diseño:} Figma, UI/UX, Principios MVC, UML

\textbf{Metodologías:} Scrum, DevSecOps, Gestión de proyectos (Jira)

\textbf{Idiomas:} Español (Nativo), Inglés (Intermedio)

% ==================== HABILIDADES INTERPERSONALES ====================
\section*{HABILIDADES INTERPERSONALES}

Resolución de problemas • Trabajo en equipo • Adaptabilidad • Gestión del tiempo • Creatividad • Comunicación efectiva • Aprendizaje continuo

\end{document}